\begin{lemma}[Smooth1LeOne]
\uses{Smooth1}
\label{Smooth1LeOne}
\lean{Smooth1LeOne}
% at /home/runner/work/trough/trough/library/PrimeNumberTheoremAnd/MellinCalculus.lean:1135.0-1190.18
\leanok
If $\nu$ is nonnegative and has mass one, then $\widetilde{1_{\epsilon}}(x)\le 1$, $\forall x>0$.
\end{lemma}
\begin{proof}
\uses{DeltaSpike,MellinConvolution}
\leanok
By Definitions \ref{Smooth1}, \ref{MellinConvolution} and \ref{DeltaSpike}
  $$
    \widetilde{1_\epsilon}(x)
    =\int_0^\infty 1_{(0,1]}(y)\frac1\epsilon\nu((x/y)^{\frac1\epsilon}) \frac{dy}y
  $$
  and since $1_{(0,1]}(y)\le 1$, and all the factors in the integrand are nonnegative,
  $$
    \widetilde{1_\epsilon}(x)\le\int_0^\infty \frac1\epsilon\nu((x/y)^{\frac1\epsilon}) \frac{dy}y
  $$
  (because in mathlib the integral of a non-integrable function is $0$, for the inequality above
  to be true, we must prove that $\nu((x/y)^{\frac1\epsilon})/y$ is integrable; this follows from
  the computation below).
  We then change variables to $z=(x/y)^{\frac1\epsilon}$:
  $$
    \widetilde{1_\epsilon}(x)\le\int_0^\infty \nu(z) \frac{dz}z
  $$
  which by Theorem \ref{SmoothExistence} is 1.
\end{proof}
